% -----------------------------------------------------------
\subsection{Arquitectura del Sistema}

La arquitectura de Bouquet se fundamenta en un modelo híbrido orientado a la
captura transaccional y al procesamiento analítico de datos. El sistema combina
una aplicación web progresiva (PWA) desplegada en Vercel, una capa transaccional
basada en PostgreSQL/Supabase y un entorno analítico batch ejecutado en un servidor
VPS mediante Apache Spark.

Esta arquitectura permite separar las responsabilidades operativas y analíticas
del sistema. La plataforma transaccional gestiona la interacción en tiempo real
entre comensales, meseros, cocina y administración, registrando eventos como
apertura de mesas virtuales, órdenes, cambios de estado y liquidaciones internas.
Posteriormente, dichos datos son extraídos y procesados en un pipeline Medallion
(Bronze, Silver y Gold) para generar indicadores de negocio y métricas de operación.

El diseño adopta un enfoque tipo Backend-for-Frontend (BFF) integrado en Next.js,
lo que permite centralizar reglas de negocio, validaciones de acceso y contratos
de comunicación hacia las distintas interfaces del sistema. De esta forma, las
vistas del comensal, mesero, cocina y administración consumen datos consistentes
desde una misma base transaccional, mientras que las consultas analíticas se
resuelven mediante tablas Gold generadas por el pipeline batch.

\subsubsection{Descripción de las Capas Arquitectónicas}

La plataforma se organiza en cinco capas principales, cada una con responsabilidades
técnicas definidas:

\begin{enumerate}[leftmargin=1.5em]
    \item \textbf{Capa de Presentación (Frontend PWA).}
    Construida con Next.js y React, permite el acceso desde navegador móvil y
    escritorio. Provee interfaces diferenciadas para comensal, mesero, cocina y
    administración. La PWA facilita una experiencia instalable sin requerir
    distribución mediante tiendas de aplicaciones.

    \item \textbf{Capa de Aplicación y API (BFF).}
    Implementada mediante rutas API, Server Actions y Middleware de Next.js. Esta
    capa valida sesiones, aplica reglas de negocio, controla permisos por rol,
    gestiona la activación de mesas virtuales, registra órdenes y coordina la
    liquidación interna. También expone los puntos de integración necesarios para
    consultar resultados analíticos desde el dashboard.

    \item \textbf{Capa de Dominio y Acceso a Datos.}
    Utiliza Prisma ORM como capa tipada de acceso a PostgreSQL. Esta capa modela
    entidades transaccionales, aplica restricciones de integridad, facilita la
    construcción de consultas seguras y permite manejar migraciones controladas
    del esquema de base de datos.

    \item \textbf{Capa de Persistencia Transaccional y Realtime.}
    Basada en Supabase/PostgreSQL. Almacena la operación OLTP del sistema: cadenas,
    zonas, sucursales, usuarios, mesas, sesiones, comensales, órdenes, estados de
    cocina, aportaciones internas y auditoría. Las políticas RLS permiten aislar datos
    por nivel organizacional. Supabase Realtime sincroniza eventos operativos como
    cambios de estado de mesa, órdenes y cocina.

    \item \textbf{Capa de Procesamiento Analítico.}
    Ejecutada en un VPS mediante Apache Spark. Extrae datos transaccionales desde
    PostgreSQL, genera datasets Bronze, aplica limpieza y enriquecimiento en Silver,
    y produce agregados Gold para el dashboard administrativo. Esta capa permite
    separar la carga analítica de la operación transaccional en tiempo real.
\end{enumerate}

\begin{table}[H]
\caption{Matriz Tecnológica y Responsabilidades por Capa}
\label{tab:capas_arq}
\centering
\small
\setstretch{1.1}
\begin{tabular}{>{\bfseries\raggedright\arraybackslash}p{3cm}
                >{\raggedright\arraybackslash}p{4cm}
                >{\raggedright\arraybackslash}p{9cm}}
\hline
\thead{Capa} & \thead{Stack Tecnológico} & \thead{Responsabilidad Primaria} \\
\hline
Presentación &
Next.js · React · Tailwind CSS &
Entrega de interfaces PWA para comensal, mesero, cocina y administración; experiencia
mobile-first y gestión de estado de interacción. \\ \rowcolor{altrow}

Aplicación / BFF &
Next.js API Routes · Server Actions · Middleware &
Validación de sesión, reglas de negocio, control de acceso, activación de mesas,
registro de órdenes y liquidación interna. \\

Dominio y Acceso a Datos &
Prisma ORM · TypeScript &
Modelado de entidades, migraciones, consultas tipadas, integridad referencial y
abstracción de persistencia. \\ \rowcolor{altrow}

Persistencia Transaccional &
Supabase · PostgreSQL · RLS · Realtime &
Almacenamiento OLTP, aislamiento multitenant vía RLS, auditoría y sincronización en tiempo
real de mesas, órdenes y cocina. \\

Procesamiento Analítico &
VPS · Apache Spark · Medallion &
Extracción batch, generación de capas Bronze/Silver/Gold, cálculo de métricas y
preparación de datos para dashboard. \\
\hline
\end{tabular}
\end{table}

A continuación, en la Figura~\ref{fig:deployment}, se presenta la representación visual de 
esta arquitectura, detallando la interacción entre los diferentes módulos y el ecosistema 
de servicios que sostienen la operación de la plataforma.

\clearpage
\begin{sidewaysfigure}
    \centering
    \includegraphics[width=\textheight]{architecture/rendered/deployment.pdf}
    \caption{Diagrama de despliegue de la arquitectura Bouquet.}
    \label{fig:deployment}
\end{sidewaysfigure}
\clearpage

